\subsection{Strict morphisms and finite-rank linear maps}

PROPOSITION 1. --- Let E and F be locally convex spaces, \(E_{1}\) a closed subspace of finite codimension of E and u a continuous linear map from E into F. For u to be a strict morphism with closed image, it is necessary and sufficient that \(u|E_{1}\) be one.

Suppose first that the linear map \(u|E_1\) is injective. We then have \(E_1 \cap \text{Ker}(u) = \{0\}\) so that \(\text{Ker}(u)\) is finite-dimensional. Let S be a vector subspace of E supplementary to \(E_1 + \text{Ker}(u)\) . The space E is the topological direct sum of \(E_1\) , \(\text{Ker}(u)\) and S (EVT, I, p. 15, cor. 4 and prop. 3). If \(u\) is a strict morphism with closed image, it defines by restriction an isomorphism from \(E_1 \oplus S\) onto a closed vector subspace of F, and a fortiori an isomorphism from \(E_1\) onto a closed vector subspace of F. Conversely, suppose that \(u\) defines by restriction an isomorphism from \(E_1\) onto \(u(E_1)\) and that \(u(E_1)\) is closed in F. We have \(u(E) = u(E_1) \oplus u(S)\) . It follows that \(u(E)\) is the topological direct sum of \(u(E_1)\) and \(u(S)\) , and is closed in F (loc. cit.). Since \(u(E_1)\) is closed and since we have \(u(E_1) \cap u(S) = \{0\}\) , the space \(u(S)\) is separated and \(u\) defines by restriction an isomorphism from S onto \(u(S)\) (EVT, I, p. 15, cor. of prop. 3), hence also of \(E_1 \oplus S\) onto \(u(E)\) . This proves that \(u\) is a strict morphism with closed image.

Moving to the general case, set \(N = E_{1} \cap \operatorname{Ker}(u)\) and G = E/N. The locally convex space \(E_{1}/N\) is identified with a closed vector subspace of finite codimension \(G_{1}\) of G. Let us denote by \(v\colon G \to F\) the continuous linear map deduced from u by passing to the quotient. For \(u\) (resp. \(u|E_{1}\)) to be a strict morphism with closed image, it is necessary and sufficient that \(v\) (resp. \(v|G_{1}\)) be one. This reduces us to the case already treated.

COROLLARY 1. --- Suppose F is separated. Let \(v \in \mathcal{L}^{\mathrm{f}}(\mathrm{E};\mathrm{F})\) . If u is a strict morphism with closed image from E into F, then the same holds for \(u + v\) .

Since F is separated, the kernel of v is a closed vector subspace of E; it has finite codimension in E and \(u + v\) has the same restriction as \(u\) on \(\operatorname{Ker}(v)\) . The corollary thus follows from the proposition.

COROLLARY 2. --- Let T be a separated finite-dimensional vector subspace of F. Denote by \(\pi\colon\mathrm{F}\to\mathrm{F}/\mathrm{T}\) the canonical surjection. For u to be a strict morphism with closed image, it is necessary and sufficient that \(\pi\circ u\) be one.

The identity map of T onto itself extends to a continuous map \(q \colon \mathrm{F} \to \mathrm{T}\) (EVT, II, p. 26, remark). The kernel S of \(q\) is a closed topological complement of T. Let \(p \colon \mathrm{F} \to \mathrm{F}\) be the projector with image S associated with the decomposition \(\mathrm{F} = \mathrm{T} \oplus \mathrm{S}\) . For \(\pi \circ u\) to be a strict morphism with closed image, it is necessary and sufficient that \(p \circ u\) be one. Moreover, \(p \circ u\) and \(u\) have the same restriction to \(u^{-1}(\mathrm{S})\) , which is a closed vector subspace of finite codimension in E, and the assertion follows from the proposition.

